


\documentclass{article}
\usepackage{amsmath, amssymb}
\usepackage{empheq}
\usepackage[a4paper, left=2.5cm, right=2.5cm, top=2cm]{geometry}
\title{Formulas Project Lutsko-Functional}
\author{Melih G\"ul  \\
	Institut f\"ur Theoretische Physik \\
	Eberhard-Karls-Universit\"at-T\"ubingen  \\
	}

\date{\today}
% Hint: \title{what ever}, \author{who care} and \date{when ever} could stand 
% before or after the \begin{document} command 
% BUT she \maketitle command MUST come AFTER the \begin{document} command! 
\begin{document}

\maketitle


\section{Rosenfeld-Functional}

In the framework of classical Density Functional Theory (cDFT) the excess free energy Functional ${\cal F}_{ex}$ for a $\nu$-component mixture is written as
\begin{equation}
	{\cal F}_{ex}[\rho_i] = \int d\mathbf{r}\,\Phi_{ex}(\{n_\alpha\}),
\end{equation}
where $\Phi_{ex}(\{n_\alpha\})$ denotes the excess free energy density as a function of weighted densities
\begin{equation}
	n_\alpha(\mathbf{r}) \equiv \sum_{i=1}^\nu\int d\mathbf{r}'\,\rho_i(\mathbf{r})\,\omega^i_\alpha(\mathbf{r}-\mathbf{r}'),
\end{equation} 
which are convolutions of the local density $\rho_i(\mathbf{r})$ and the weight functions $\omega^i_\alpha(\mathbf{r}-\mathbf{r}')$ of species $i$ given in FMT.\\
Employing thermodynamic relations and scaled particle theory (SPT) one can derive a differential equation for $\Phi_{ex}$. It reads
\begin{equation}\label{DiffPhi}
	\frac{\partial\Phi_{ex}}{\partial n_3} = n_0-\Phi_{ex}+\sum_\alpha\frac{\partial\Phi_{ex}}{\partial n_\alpha}n_\alpha,
\end{equation}
where the left-hand side of the equation is related to the bulk-pressure $\beta p$ of the system.
For the 1D hard-sphere system, Eq.\eqref{DiffPhi} yields the exact solution 
\begin{equation}
	\Phi_{ex}^{(1D)} = -n_0\log(1-n_1).
\end{equation}
For higher dimensions however, only solutions with the right low-density limit are known.\\ 
The Rosenfeld-Functional as a solution of Eq.\eqref{DiffPhi} in the 3D case is given by the excess free energy density 
\begin{equation}
	\Phi^{\text{RF}}_{ex}(\{n_\alpha\}) = \Phi_{ex}^{(1)}+\Phi_{ex}^{(2)}+\Phi_{ex}^{(3)},
\end{equation}
where
\begin{align}\label{Rosenfeld}
	\Phi_{ex}^{(1)} &= -n_0\log(1-n_3),\quad
	\Phi_{ex}^{(2)} = \frac{n_1 n_2-\mathbf{n}_1 \mathbf{n}_2}{1-n_3},\quad
	\Phi_{ex}^{(3)} = \frac{n_2^3 - 3n_2\mathbf{n}_2^2}{24\pi (1-n_3)^2}.
\end{align}
Further improvements with respect to the functional were achieved by using different equation of states (EOS) as an input of Eq.\eqref{DiffPhi}. Taking the Carnahan-Starling EOS for instance leads to the White-Bear solution where
\begin{equation}
	\Phi_{ex}^{\text{WB},(3)}(\{n_\alpha\}) = \frac{n_3+(1-n_3)^2\log(1-n_3)}{36\pi n_3^2(1-n_3)^2}(n_2^3 - 3n_2\mathbf{n}_2^2)
\end{equation}
and the other parts remaining unchanged. A Mark-II version of the WB-functional could be found using a generalized version of the CS-EOS for hard-sphere mixtures.

\section{Lutsko-Functional}

Lutsko presented a formalism based on dimensional crossover that is capable of reproducing the correct low-dimensional version given a functional in higher dimension.\\
He could show that $\Phi_{ex}$ takes the following form
\begin{equation}
	\Phi_{ex}^{\text{L}}(\{n_\alpha\}) = \Phi_{ex}^{(1)}(\{n_\alpha\})+\Phi_{ex}^{(2)}(\{n_\alpha\})+\Phi_{ex}^{(3)}(\{n_\alpha\};A,B),
\end{equation}
where $\Phi_{ex}^{(1)}(\{n_\alpha\})$ and $\Phi_{ex}^{(2)}(\{n_\alpha\})$ equal those of Eq.\eqref{Rosenfeld}. $\Phi_{ex}^{(3)}(\{n_\alpha\};A,B)$ now includes two free parameters $A$ and $B$
\begin{equation}
	\Phi_{ex}^{(3)}(\{n_\alpha\};A,B) = \frac{(A+B)n_2^3-3An_2\mathbf{n}_2^2+3A\mathbf{n}_2\mathbf{T}\mathbf{n}_2-3Bn_2\text{Tr}\,\mathbf{T}^2+(2B-A)\text{Tr}\,\mathbf{T}^3}{24\pi(1-n_3)^2}.
\end{equation}
Apparently, this functional's performance depends on the choice being made for $A$ and $B$. Imposing certain conditions could fix their value. 

\section{Sum Rules}

The conditions we employ make use of the sum rules with respect to excess chemical potential $\mu_{ex}$ and isothermal compressibility $\chi_T$. Both can be obtained in two ways and by minimizing the difference, i.e. looking for the best agreement between these two approaches, will fix these parameters.\\
$\mu_{ex}$ can be calculated first by evaluating the grand potential $\Omega$ which is obtained from a minimization program and second from SPT. $\chi_T$ is obtained on the one hand by making use of the thermodynamic derivation and on the other hand calculating the structure factor $S(0)$.\\
Hence, demanding consistency provides two constraints on two free parameters.
It is also possible to investigate the direct correlation function and try to optimize in this way.\\   
In total, we have
\begin{itemize}
	\item Thermodynamics\\
		
		\begin{equation}
			\beta\mu_{ex} = v\frac{\partial f_{ex}}{\partial\eta},\quad \chi_T = \frac{1}{\eta}\left(\frac{\partial\beta\mu}{\partial\eta}\right)_{T,V}^{-1}.
		\end{equation}
		
	\item DFT\\
	
		\begin{equation}
			\beta\mu_{ex} = \beta pv + \beta\gamma A,\quad \chi_T = 1 + 4\pi\rho\int_0^{\infty}dr\,r^2h(r).
		\end{equation}
\end{itemize}
Here, $v$ is the particle volume, $\rho$ the bulk density and $\eta = \rho v$ the packing fraction. The pair direct correlation function $h(r)$ can be easily obtained from the density profile $\rho(r)$ being derived from a minimization program. We also assumed spherical symmetry.

\section{Direct Correlation Function (DCF)}

In FMT, the DCF is defined as
\begin{equation}
	c^{(2)}(\mathbf{r}-\mathbf{r}') \equiv -\frac{\delta^2\beta{\cal F}_{ex}[\rho]}{\delta\rho(\mathbf{r})\delta\rho(\mathbf{r}')}\Bigg|_{\rho(\mathbf{r})=\rho_b,\,\rho(\mathbf{r}')=\rho_b},
\end{equation}
which after applying the chain rule becomes
\begin{equation}
	c^{(2)}(\mathbf{r}-\mathbf{r}') = -\sum_{\alpha,\beta}\frac{\partial^2\Phi_{ex}}{\partial n_\alpha\partial n_\beta}\,\omega_\alpha\otimes\omega_\beta(\mathbf{r}-\mathbf{r}'),
\end{equation}
where $\omega_\alpha\otimes\omega_\beta$ denotes
\begin{equation}
	\omega_\alpha\otimes\omega_\beta(\mathbf{r}-\mathbf{r}')= \int d\mathbf{r}''\,\omega_\alpha(\mathbf{r}-\mathbf{r}'')\omega_\beta(\mathbf{r}'-\mathbf{r}'').
\end{equation}


\end{document}